\documentclass[english,11pt,a4paper]{article}
\usepackage[margin=1in]{geometry}
\usepackage{fontspec}
\setmainfont{lmroman10-regular.otf}[BoldFont=lmroman10-bold.otf,ItalicFont=lmroman10-italic.otf,BoldItalicFont=lmroman10-bolditalic.otf]
\setsansfont{lmsans10-regular.otf}[BoldFont=lmsans10-bold.otf]
\setmonofont{lmmono10-regular.otf}[Scale=MatchLowercase]
\newfontfamily\LeanSymbols{seguisym.ttf}[Path=typeset_runtime/]
\newlength{\LeanCharWidth}
\usepackage{amsmath,amssymb}
\usepackage[english]{babel}
\usepackage{microtype,parskip}
\usepackage{graphicx,float,caption}
\usepackage{longtable,booktabs,array,calc}
\usepackage{xcolor}
\usepackage{xurl}
\usepackage{hyperref,bookmark}
\usepackage{etoolbox,needspace}
\setcounter{secnumdepth}{-\maxdimen}
\setlength{\emergencystretch}{3em}
\providecommand{\tightlist}{\setlength{\itemsep}{0pt}\setlength{\parskip}{0pt}}
\newcounter{none}
\AtBeginEnvironment{longtable}{\Needspace{8\baselineskip}\small}
\pretocmd{\section}{\Needspace{5\baselineskip}}{}{}
\pretocmd{\subsection}{\Needspace{4\baselineskip}}{}{}
\pretocmd{\subsubsection}{\Needspace{4\baselineskip}}{}{}
\captionsetup{font=small,labelfont=bf}
\urlstyle{same}
\providecommand{\pandocbounded}[1]{#1}
\hypersetup{pdftitle={Formal companion to Paper IV},pdfauthor={Juan Pablo Traverso Gianini},pdflang={en-US},colorlinks=true,linkcolor=blue,urlcolor=blue,citecolor=blue}
\title{Formal companion to Paper IV}
\author{Juan Pablo Traverso Gianini\\\small Independent researcher, Santiago, Chile\\\small \href{mailto:jtraverso@gmail.com}{jtraverso@gmail.com}\\\small ORCID: 0009-0003-6068-4096}
\date{}
\begin{document}
\raggedbottom
\maketitle
\textbf{Preprint · version 1.24}\\
October 19, 2026

This document describes the frozen formalization supporting \emph{Clique
partitions with rooted simplicial defect: quantitative stability and
sharp eventual bounds}. It accompanies the editorial revision of the
manuscript; the mathematical statements and the Lean source freeze are
unchanged. Section, equation and theorem numbers refer to the main
manuscript unless otherwise specified. References {[}1--23{]} retain its
bibliography.

\subsection{1. Source identity and verification
scope}\label{source-identity-and-verification-scope}

{\def\LTcaptype{none} % do not increment counter
\begin{longtable}[]{@{}ll@{}}
\toprule\noalign{}
Item & Frozen value \\
\midrule\noalign{}
\endhead
\bottomrule\noalign{}
\endlastfoot
Lean and Mathlib & v4.28.0 \\
Source cut & \nolinkurl{piv-v12-fb459343d234} \\
Main source modules & 607 \\
Source/configuration/documentation manifest entries & 615 \\
Selected audit targets & 19 \\
Literal exports through \texttt{import\ PaperIV} & 224 \\
Recorded axiom outputs & 461; not a count of distinct theorems \\
\end{longtable}
}

The source manifest SHA-256 is
\nolinkurl{fb459343d234f968d7d32eff1491ea8a09aa2e135b313a80623012e7449042f5}.
The source-only archive SHA-256 is
\nolinkurl{cb2741454380736ffe01b59933057b5079fa9f865d5bbdd3f67534bf808a62e6}.
These identify the original formal source package, not the new editorial
PDFs.

The
\href{https://github.com/jtraverso/erdos-81-chordal-clique-partitions/tree/ef31294ad67dfe3f44bc47cec4791c836f55e8aa/preprints/PAPER_IV/05_formalization}{frozen
package} and
\href{https://github.com/jtraverso/erdos-81-chordal-clique-partitions/tree/ef31294ad67dfe3f44bc47cec4791c836f55e8aa/preprints/PAPER_IV/03_reproducibility}{recorded
build evidence} are available in the series repository. The series
concept DOI is
\href{https://doi.org/10.5281/zenodo.21273143}{10.5281/zenodo.21273143}.
The preceding manuscript edition and its deposit remain available; the
source freeze is the same.

The identified Lean supplement for this version uses Lean and Mathlib
v4.28.0. The final declarations do not take rounding, existence of an
optimum or the near constructor as arguments: they invoke them as proved
theorems in their dependencies. The permitted foundational axioms are
distinguished from any additional mathematical hypothesis.

The recorded reconstruction covers all 607 modules: 421 objects built in
the first segment were verified and reused after a restart, and the
resumed segment compiled 186 modules. All 19 selected targets were
freshly executed with exit code zero. Their 461 axiom outputs use only
\nolinkurl{propext}, \nolinkurl{Classical.choice} and \nolinkurl{Quot.sound}, or
subsets. The records distinguish source identity, compilation, export
availability and transitive axiom checks. None of these counts is a
claim that every manuscript sentence is a separately exported theorem.

The consolidated external report for v1.22-r4 records PASS for E0--E8
within its documented scope. That report and the build logs remain
evidence for the identified frozen source. They are not a fresh external
review of v1.24. The present editorial review checks relocation,
statement preservation and document rendering; no Lean build or new
mathematical verification is claimed.

\subsection{2. Mathematical statements and exported
declarations}\label{mathematical-statements-and-exported-declarations}

The common prefix \nolinkurl{PaperIV} is omitted where indicated. The
following map preserves the original theorem labels; moving it out of
the main manuscript has not changed the formal interface.

{\def\LTcaptype{none} % do not increment counter
\begin{longtable}[]{@{}
  >{\raggedright\arraybackslash}p{(\linewidth - 2\tabcolsep) * \real{0.5000}}
  >{\raggedright\arraybackslash}p{(\linewidth - 2\tabcolsep) * \real{0.5000}}@{}}
\toprule\noalign{}
\begin{minipage}[b]{\linewidth}\raggedright
Result
\end{minipage} & \begin{minipage}[b]{\linewidth}\raggedright
Lean declaration
\end{minipage} \\
\midrule\noalign{}
\endhead
\bottomrule\noalign{}
\endlastfoot
Theorem A, additive form &
\nolinkurl{Erdos81AllOrders.erdos81_all_orders_additive} \\
Theorem B, upper bound &
\nolinkurl{Erdos81Unconditional.erdos81_cliquePartition} \\
Theorem B, maximum & \nolinkurl{PaperTheorems.erdos81_max_eq} \\
Theorem 3.1 & \nolinkurl{RC01Final.rc01_uniformRoundingTarget} \\
Corollary 3.5 &
\nolinkurl{FarRegimeAllGraphs.farRegime_cliquePartition_allGraphs} \\
Dichotomy (1.5) &
\nolinkurl{HybridDichotomy.chordal_far_or_nearStructure} \\
Theorem 5.0, strengthened window &
\nolinkurl{NearH1LocalConstructor.exists_near_partition_paid_by_root_sharp} \\
Theorem 5.0, previous interface &
\nolinkurl{NearH1LocalConstructor.exists_near_partition_paid_by_root} \\
Corollary 5.4, size and density &
\nolinkurl{NearCriticalDichotomy.chordal_far_or_criticalRoot},
\nolinkurl{chordal_far_or_edge_density} \\
Optimal quadratic coefficient &
\nolinkurl{SharpConstantOptimality.erdos81_quadratic_constant_optimal},
\nolinkurl{erdos81_quadratic_constant_isLeast} \\
Theorem 6.1, strengthened form &
\nolinkurl{IntegralStability.chordal_linear_stability_sixteen} \\
Theorem 6.1, earlier interface &
\nolinkurl{IntegralStability.chordal_linear_stability} \\
Corollary 6.1a, piece count &
\nolinkurl{RootPartitionStability.chordal_partition_stability} \\
Corollary 6.1a, joint real piece and edge bounds &
\nolinkurl{SublinearResearch.chordal_joint_stability_real} in
\nolinkurl{PublicationStability} \\
Identity (6.7b) &
\nolinkurl{RootPartitionStability.sum_rootPieceDefect_eq} \\
Corollary 6.2 &
\nolinkurl{ExtremalClassification.chordal_extremal_classification} \\
Proposition D.3 & \nolinkurl{SplitMixedGap.mixed_gap_zero} together with
the bounded construction in \nolinkurl{SplitCompleteSharpValue} \\
\end{longtable}
}

For the extension table below, the prefix is
\nolinkurl{PaperIV.SublinearResearch}, except for the named E32 entry.

{\def\LTcaptype{none} % do not increment counter
\begin{longtable}[]{@{}
  >{\raggedright\arraybackslash}p{(\linewidth - 2\tabcolsep) * \real{0.5000}}
  >{\raggedright\arraybackslash}p{(\linewidth - 2\tabcolsep) * \real{0.5000}}@{}}
\toprule\noalign{}
\begin{minipage}[b]{\linewidth}\raggedright
Statement
\end{minipage} & \begin{minipage}[b]{\linewidth}\raggedright
Declaration
\end{minipage} \\
\midrule\noalign{}
\endhead
\bottomrule\noalign{}
\endlastfoot
Real same-root chordal bounds, including edge mass &
\nolinkurl{chordal_joint_stability_real} \\
Theorem \(C^\prime\) and (6.12)--(6.13) &
\nolinkurl{fixed_defect_joint_stability_real} \\
Explicit parameters and classification &
\nolinkurl{FixedExplicit.fixed_defect_stability_explicit} \\
Optimal-size edit bound (6.17) &
\nolinkurl{fixed_defect_exact_extremal_edit_real} \\
Unrestricted equality classification &
\nolinkurl{E32.cp_classification_of_theoremCPrimeC}, with the proved
classification above and the historical induced-removal assembly
{[}23{]} (see companion §3) \\
Finite bound (6.22) &
\nolinkurl{certified_shifted_fractional_bound_all_orders} \\
Finite bound (6.23) &
\nolinkurl{certified_refined_fractional_bound_all_orders} \\
Uniform integral versions & \nolinkurl{uniform_shifted_partition_bound};
\nolinkurl{uniform_refined_partition_bound} \\
Arbitrary-order sequence approximation &
\nolinkurl{sequence_arbitrary_orders_partition_bound} \\
One-root sequence stability &
\nolinkurl{sequence_arbitrary_orders_stability} \\
\end{longtable}
}

The finite same-root interfaces take real deficits. The sequence
interface records rational error and excess parameters. That difference
is retained when comparing the manuscript with elaborated types.

\textbf{Literal definitions.} The following excerpts reproduce
definitions and theorem headers, omitting proofs but not hypotheses. The
first excerpt is in the namespace \nolinkurl{SimpleGraph}; the next three
are in \nolinkurl{PaperIV.FarRounding}, with finite \(V\), decidable
equality and decidable adjacency for \(G\). The last definition is in
\nolinkurl{PaperIV}. The source locations are
\nolinkurl{PaperIV/ChordalStructure.lean} (chordality),
\nolinkurl{PaperIV/FarRounding.lean} (pairs and partitions), and
\nolinkurl{PaperIV/TargetEnvelope.lean} (target size), within the frozen
package linked in §1 of this companion.

\begin{quote}\fontsize{8.5}{10.2}\selectfont\ttfamily\raggedright\setlength{\parskip}{0pt}\settowidth{\LeanCharWidth}{0}
\noindent def\hspace*{\dimexpr 1\LeanCharWidth\relax}IsChordal\hspace*{\dimexpr 1\LeanCharWidth\relax}(G\hspace*{\dimexpr 1\LeanCharWidth\relax}:\hspace*{\dimexpr 1\LeanCharWidth\relax}SimpleGraph\hspace*{\dimexpr 1\LeanCharWidth\relax}V)\hspace*{\dimexpr 1\LeanCharWidth\relax}:\hspace*{\dimexpr 1\LeanCharWidth\relax}Prop\hspace*{\dimexpr 1\LeanCharWidth\relax}:=\par
\noindent \hspace*{\dimexpr 2\LeanCharWidth\relax}{\LeanSymbols ∀}\hspace*{\dimexpr 1\LeanCharWidth\relax}{\LeanSymbols ⦃}v\hspace*{\dimexpr 1\LeanCharWidth\relax}:\hspace*{\dimexpr 1\LeanCharWidth\relax}V{\LeanSymbols ⦄}\hspace*{\dimexpr 1\LeanCharWidth\relax}(c\hspace*{\dimexpr 1\LeanCharWidth\relax}:\hspace*{\dimexpr 1\LeanCharWidth\relax}G.Walk\hspace*{\dimexpr 1\LeanCharWidth\relax}v\hspace*{\dimexpr 1\LeanCharWidth\relax}v),\hspace*{\dimexpr 1\LeanCharWidth\relax}c.IsCycle\hspace*{\dimexpr 1\LeanCharWidth\relax}{\LeanSymbols →}\hspace*{\dimexpr 1\LeanCharWidth\relax}4\hspace*{\dimexpr 1\LeanCharWidth\relax}{\LeanSymbols ≤}\hspace*{\dimexpr 1\LeanCharWidth\relax}c.length\hspace*{\dimexpr 1\LeanCharWidth\relax}{\LeanSymbols →}\par
\noindent \hspace*{\dimexpr 4\LeanCharWidth\relax}{\LeanSymbols ∃}\hspace*{\dimexpr 1\LeanCharWidth\relax}x\hspace*{\dimexpr 1\LeanCharWidth\relax}y\hspace*{\dimexpr 1\LeanCharWidth\relax}:\hspace*{\dimexpr 1\LeanCharWidth\relax}V,\hspace*{\dimexpr 1\LeanCharWidth\relax}x\hspace*{\dimexpr 1\LeanCharWidth\relax}{\LeanSymbols ∈}\hspace*{\dimexpr 1\LeanCharWidth\relax}c.support\hspace*{\dimexpr 1\LeanCharWidth\relax}{\LeanSymbols ∧}\hspace*{\dimexpr 1\LeanCharWidth\relax}y\hspace*{\dimexpr 1\LeanCharWidth\relax}{\LeanSymbols ∈}\hspace*{\dimexpr 1\LeanCharWidth\relax}c.support\hspace*{\dimexpr 1\LeanCharWidth\relax}{\LeanSymbols ∧}\hspace*{\dimexpr 1\LeanCharWidth\relax}G.Adj\hspace*{\dimexpr 1\LeanCharWidth\relax}x\hspace*{\dimexpr 1\LeanCharWidth\relax}y\hspace*{\dimexpr 1\LeanCharWidth\relax}{\LeanSymbols ∧}\hspace*{\dimexpr 1\LeanCharWidth\relax}s(x,\hspace*{\dimexpr 1\LeanCharWidth\relax}y)\hspace*{\dimexpr 1\LeanCharWidth\relax}{\LeanSymbols ∉}\hspace*{\dimexpr 1\LeanCharWidth\relax}c.edges\par
\end{quote}

\begin{quote}\fontsize{8.5}{10.2}\selectfont\ttfamily\raggedright\setlength{\parskip}{0pt}\settowidth{\LeanCharWidth}{0}
\noindent def\hspace*{\dimexpr 1\LeanCharWidth\relax}pairs\hspace*{\dimexpr 1\LeanCharWidth\relax}(K\hspace*{\dimexpr 1\LeanCharWidth\relax}:\hspace*{\dimexpr 1\LeanCharWidth\relax}Finset\hspace*{\dimexpr 1\LeanCharWidth\relax}V)\hspace*{\dimexpr 1\LeanCharWidth\relax}:\hspace*{\dimexpr 1\LeanCharWidth\relax}Finset\hspace*{\dimexpr 1\LeanCharWidth\relax}(Sym2\hspace*{\dimexpr 1\LeanCharWidth\relax}V)\hspace*{\dimexpr 1\LeanCharWidth\relax}:=\hspace*{\dimexpr 1\LeanCharWidth\relax}K.sym2.filter\hspace*{\dimexpr 1\LeanCharWidth\relax}fun\hspace*{\dimexpr 1\LeanCharWidth\relax}e\hspace*{\dimexpr 1\LeanCharWidth\relax}=>\hspace*{\dimexpr 1\LeanCharWidth\relax}{\LeanSymbols ¬}\hspace*{\dimexpr 1\LeanCharWidth\relax}e.IsDiag\par
\noindent \strut\par
\noindent structure\hspace*{\dimexpr 1\LeanCharWidth\relax}CliquePartition\hspace*{\dimexpr 1\LeanCharWidth\relax}where\par
\noindent \hspace*{\dimexpr 2\LeanCharWidth\relax}pieces\hspace*{\dimexpr 1\LeanCharWidth\relax}:\hspace*{\dimexpr 1\LeanCharWidth\relax}Finset\hspace*{\dimexpr 1\LeanCharWidth\relax}(Finset\hspace*{\dimexpr 1\LeanCharWidth\relax}V)\par
\noindent \hspace*{\dimexpr 2\LeanCharWidth\relax}isClique\hspace*{\dimexpr 1\LeanCharWidth\relax}:\hspace*{\dimexpr 1\LeanCharWidth\relax}{\LeanSymbols ∀}\hspace*{\dimexpr 1\LeanCharWidth\relax}K\hspace*{\dimexpr 1\LeanCharWidth\relax}{\LeanSymbols ∈}\hspace*{\dimexpr 1\LeanCharWidth\relax}pieces,\hspace*{\dimexpr 1\LeanCharWidth\relax}{\LeanSymbols ∀}\hspace*{\dimexpr 1\LeanCharWidth\relax}a\hspace*{\dimexpr 1\LeanCharWidth\relax}{\LeanSymbols ∈}\hspace*{\dimexpr 1\LeanCharWidth\relax}K,\hspace*{\dimexpr 1\LeanCharWidth\relax}{\LeanSymbols ∀}\hspace*{\dimexpr 1\LeanCharWidth\relax}b\hspace*{\dimexpr 1\LeanCharWidth\relax}{\LeanSymbols ∈}\hspace*{\dimexpr 1\LeanCharWidth\relax}K,\hspace*{\dimexpr 1\LeanCharWidth\relax}a\hspace*{\dimexpr 1\LeanCharWidth\relax}{\LeanSymbols ≠}\hspace*{\dimexpr 1\LeanCharWidth\relax}b\hspace*{\dimexpr 1\LeanCharWidth\relax}{\LeanSymbols →}\hspace*{\dimexpr 1\LeanCharWidth\relax}G.Adj\hspace*{\dimexpr 1\LeanCharWidth\relax}a\hspace*{\dimexpr 1\LeanCharWidth\relax}b\par
\noindent \hspace*{\dimexpr 2\LeanCharWidth\relax}two\_le\_card\hspace*{\dimexpr 1\LeanCharWidth\relax}:\hspace*{\dimexpr 1\LeanCharWidth\relax}{\LeanSymbols ∀}\hspace*{\dimexpr 1\LeanCharWidth\relax}K\hspace*{\dimexpr 1\LeanCharWidth\relax}{\LeanSymbols ∈}\hspace*{\dimexpr 1\LeanCharWidth\relax}pieces,\hspace*{\dimexpr 1\LeanCharWidth\relax}2\hspace*{\dimexpr 1\LeanCharWidth\relax}{\LeanSymbols ≤}\hspace*{\dimexpr 1\LeanCharWidth\relax}K.card\par
\noindent \hspace*{\dimexpr 2\LeanCharWidth\relax}edgeDisjoint\hspace*{\dimexpr 1\LeanCharWidth\relax}:\hspace*{\dimexpr 1\LeanCharWidth\relax}{\LeanSymbols ∀}\hspace*{\dimexpr 1\LeanCharWidth\relax}K\hspace*{\dimexpr 1\LeanCharWidth\relax}{\LeanSymbols ∈}\hspace*{\dimexpr 1\LeanCharWidth\relax}pieces,\hspace*{\dimexpr 1\LeanCharWidth\relax}{\LeanSymbols ∀}\hspace*{\dimexpr 1\LeanCharWidth\relax}L\hspace*{\dimexpr 1\LeanCharWidth\relax}{\LeanSymbols ∈}\hspace*{\dimexpr 1\LeanCharWidth\relax}pieces,\hspace*{\dimexpr 1\LeanCharWidth\relax}K\hspace*{\dimexpr 1\LeanCharWidth\relax}{\LeanSymbols ≠}\hspace*{\dimexpr 1\LeanCharWidth\relax}L\hspace*{\dimexpr 1\LeanCharWidth\relax}{\LeanSymbols →}\hspace*{\dimexpr 1\LeanCharWidth\relax}Disjoint\hspace*{\dimexpr 1\LeanCharWidth\relax}(pairs\hspace*{\dimexpr 1\LeanCharWidth\relax}K)\hspace*{\dimexpr 1\LeanCharWidth\relax}(pairs\hspace*{\dimexpr 1\LeanCharWidth\relax}L)\par
\noindent \hspace*{\dimexpr 2\LeanCharWidth\relax}covers\hspace*{\dimexpr 1\LeanCharWidth\relax}:\hspace*{\dimexpr 1\LeanCharWidth\relax}pieces.biUnion\hspace*{\dimexpr 1\LeanCharWidth\relax}pairs\hspace*{\dimexpr 1\LeanCharWidth\relax}=\hspace*{\dimexpr 1\LeanCharWidth\relax}G.edgeFinset\par
\noindent \strut\par
\noindent def\hspace*{\dimexpr 1\LeanCharWidth\relax}CliquePartition.size\hspace*{\dimexpr 1\LeanCharWidth\relax}(Q\hspace*{\dimexpr 1\LeanCharWidth\relax}:\hspace*{\dimexpr 1\LeanCharWidth\relax}CliquePartition\hspace*{\dimexpr 1\LeanCharWidth\relax}G)\hspace*{\dimexpr 1\LeanCharWidth\relax}:\hspace*{\dimexpr 1\LeanCharWidth\relax}{\LeanSymbols ℕ}\hspace*{\dimexpr 1\LeanCharWidth\relax}:=\hspace*{\dimexpr 1\LeanCharWidth\relax}Q.pieces.card\par
\noindent \strut\par
\noindent def\hspace*{\dimexpr 1\LeanCharWidth\relax}CliquePartition.OrderAtMost\hspace*{\dimexpr 1\LeanCharWidth\relax}(Q\hspace*{\dimexpr 1\LeanCharWidth\relax}:\hspace*{\dimexpr 1\LeanCharWidth\relax}CliquePartition\hspace*{\dimexpr 1\LeanCharWidth\relax}G)\hspace*{\dimexpr 1\LeanCharWidth\relax}(r\hspace*{\dimexpr 1\LeanCharWidth\relax}:\hspace*{\dimexpr 1\LeanCharWidth\relax}{\LeanSymbols ℕ})\hspace*{\dimexpr 1\LeanCharWidth\relax}:\hspace*{\dimexpr 1\LeanCharWidth\relax}Prop\hspace*{\dimexpr 1\LeanCharWidth\relax}:=\par
\noindent \hspace*{\dimexpr 2\LeanCharWidth\relax}{\LeanSymbols ∀}\hspace*{\dimexpr 1\LeanCharWidth\relax}K\hspace*{\dimexpr 1\LeanCharWidth\relax}{\LeanSymbols ∈}\hspace*{\dimexpr 1\LeanCharWidth\relax}Q.pieces,\hspace*{\dimexpr 1\LeanCharWidth\relax}K.card\hspace*{\dimexpr 1\LeanCharWidth\relax}{\LeanSymbols ≤}\hspace*{\dimexpr 1\LeanCharWidth\relax}r\par
\end{quote}

\begin{quote}\fontsize{8.5}{10.2}\selectfont\ttfamily\raggedright\setlength{\parskip}{0pt}\settowidth{\LeanCharWidth}{0}
\noindent def\hspace*{\dimexpr 1\LeanCharWidth\relax}targetSize\hspace*{\dimexpr 1\LeanCharWidth\relax}(n\hspace*{\dimexpr 1\LeanCharWidth\relax}:\hspace*{\dimexpr 1\LeanCharWidth\relax}{\LeanSymbols ℕ})\hspace*{\dimexpr 1\LeanCharWidth\relax}:\hspace*{\dimexpr 1\LeanCharWidth\relax}{\LeanSymbols ℕ}\hspace*{\dimexpr 1\LeanCharWidth\relax}:=\hspace*{\dimexpr 1\LeanCharWidth\relax}n\hspace*{\dimexpr 1\LeanCharWidth\relax}*\hspace*{\dimexpr 1\LeanCharWidth\relax}(n\hspace*{\dimexpr 1\LeanCharWidth\relax}+\hspace*{\dimexpr 1\LeanCharWidth\relax}1)\hspace*{\dimexpr 1\LeanCharWidth\relax}/\hspace*{\dimexpr 1\LeanCharWidth\relax}6\par
\end{quote}

Natural-number division implements the floor in \(M(n)\). The expository
symbol \(c_4\) denotes the minimum over these partitions with
\texttt{OrderAtMost\ 4}; the exported statement does not assume a
minimum function, but directly supplies a witness partition. Chordality
in \nolinkurl{FarRounding} is a literal alias of the preceding cycle-based
definition.

\textbf{Exported theorems.} In \nolinkurl{PaperIV.Erdos81AllOrders} and
\nolinkurl{PaperIV.Erdos81Unconditional}, respectively, the exact headers
are:

\begin{quote}\fontsize{8.5}{10.2}\selectfont\ttfamily\raggedright\setlength{\parskip}{0pt}\settowidth{\LeanCharWidth}{0}
\noindent theorem\hspace*{\dimexpr 1\LeanCharWidth\relax}erdos81\_all\_orders\_additive\hspace*{\dimexpr 1\LeanCharWidth\relax}:\par
\noindent \hspace*{\dimexpr 4\LeanCharWidth\relax}{\LeanSymbols ∃}\hspace*{\dimexpr 1\LeanCharWidth\relax}b\hspace*{\dimexpr 1\LeanCharWidth\relax}:\hspace*{\dimexpr 1\LeanCharWidth\relax}{\LeanSymbols ℕ},\hspace*{\dimexpr 1\LeanCharWidth\relax}{\LeanSymbols ∀}\hspace*{\dimexpr 1\LeanCharWidth\relax}(n\hspace*{\dimexpr 1\LeanCharWidth\relax}:\hspace*{\dimexpr 1\LeanCharWidth\relax}{\LeanSymbols ℕ})\hspace*{\dimexpr 1\LeanCharWidth\relax}(G\hspace*{\dimexpr 1\LeanCharWidth\relax}:\hspace*{\dimexpr 1\LeanCharWidth\relax}SimpleGraph\hspace*{\dimexpr 1\LeanCharWidth\relax}(Fin\hspace*{\dimexpr 1\LeanCharWidth\relax}n))\hspace*{\dimexpr 1\LeanCharWidth\relax}[DecidableRel\hspace*{\dimexpr 1\LeanCharWidth\relax}G.Adj],\par
\noindent \hspace*{\dimexpr 6\LeanCharWidth\relax}PaperIV.FarRounding.IsChordal\hspace*{\dimexpr 1\LeanCharWidth\relax}G\hspace*{\dimexpr 1\LeanCharWidth\relax}{\LeanSymbols →}\par
\noindent \hspace*{\dimexpr 8\LeanCharWidth\relax}{\LeanSymbols ∃}\hspace*{\dimexpr 1\LeanCharWidth\relax}Q\hspace*{\dimexpr 1\LeanCharWidth\relax}:\hspace*{\dimexpr 1\LeanCharWidth\relax}CliquePartition\hspace*{\dimexpr 1\LeanCharWidth\relax}G,\hspace*{\dimexpr 1\LeanCharWidth\relax}Q.OrderAtMost\hspace*{\dimexpr 1\LeanCharWidth\relax}4\hspace*{\dimexpr 1\LeanCharWidth\relax}{\LeanSymbols ∧}\hspace*{\dimexpr 1\LeanCharWidth\relax}Q.size\hspace*{\dimexpr 1\LeanCharWidth\relax}{\LeanSymbols ≤}\hspace*{\dimexpr 1\LeanCharWidth\relax}PaperIV.targetSize\hspace*{\dimexpr 1\LeanCharWidth\relax}n\hspace*{\dimexpr 1\LeanCharWidth\relax}+\hspace*{\dimexpr 1\LeanCharWidth\relax}b\par
\end{quote}

\begin{quote}\fontsize{8.5}{10.2}\selectfont\ttfamily\raggedright\setlength{\parskip}{0pt}\settowidth{\LeanCharWidth}{0}
\noindent theorem\hspace*{\dimexpr 1\LeanCharWidth\relax}erdos81\_cliquePartition\hspace*{\dimexpr 1\LeanCharWidth\relax}:\par
\noindent \hspace*{\dimexpr 4\LeanCharWidth\relax}{\LeanSymbols ∃}\hspace*{\dimexpr 1\LeanCharWidth\relax}N\hspace*{\dimexpr 1\LeanCharWidth\relax}:\hspace*{\dimexpr 1\LeanCharWidth\relax}{\LeanSymbols ℕ},\hspace*{\dimexpr 1\LeanCharWidth\relax}{\LeanSymbols ∀}\hspace*{\dimexpr 1\LeanCharWidth\relax}n\hspace*{\dimexpr 1\LeanCharWidth\relax}:\hspace*{\dimexpr 1\LeanCharWidth\relax}{\LeanSymbols ℕ},\hspace*{\dimexpr 1\LeanCharWidth\relax}N\hspace*{\dimexpr 1\LeanCharWidth\relax}{\LeanSymbols ≤}\hspace*{\dimexpr 1\LeanCharWidth\relax}n\hspace*{\dimexpr 1\LeanCharWidth\relax}{\LeanSymbols →}\hspace*{\dimexpr 1\LeanCharWidth\relax}{\LeanSymbols ∀}\hspace*{\dimexpr 1\LeanCharWidth\relax}(G\hspace*{\dimexpr 1\LeanCharWidth\relax}:\hspace*{\dimexpr 1\LeanCharWidth\relax}SimpleGraph\hspace*{\dimexpr 1\LeanCharWidth\relax}(Fin\hspace*{\dimexpr 1\LeanCharWidth\relax}n))\hspace*{\dimexpr 1\LeanCharWidth\relax}[DecidableRel\hspace*{\dimexpr 1\LeanCharWidth\relax}G.Adj],\par
\noindent \hspace*{\dimexpr 6\LeanCharWidth\relax}PaperIV.FarRounding.IsChordal\hspace*{\dimexpr 1\LeanCharWidth\relax}G\hspace*{\dimexpr 1\LeanCharWidth\relax}{\LeanSymbols →}\par
\noindent \hspace*{\dimexpr 8\LeanCharWidth\relax}{\LeanSymbols ∃}\hspace*{\dimexpr 1\LeanCharWidth\relax}Q\hspace*{\dimexpr 1\LeanCharWidth\relax}:\hspace*{\dimexpr 1\LeanCharWidth\relax}CliquePartition\hspace*{\dimexpr 1\LeanCharWidth\relax}G,\hspace*{\dimexpr 1\LeanCharWidth\relax}Q.OrderAtMost\hspace*{\dimexpr 1\LeanCharWidth\relax}4\hspace*{\dimexpr 1\LeanCharWidth\relax}{\LeanSymbols ∧}\hspace*{\dimexpr 1\LeanCharWidth\relax}Q.size\hspace*{\dimexpr 1\LeanCharWidth\relax}{\LeanSymbols ≤}\hspace*{\dimexpr 1\LeanCharWidth\relax}targetSize\hspace*{\dimexpr 1\LeanCharWidth\relax}n\par
\end{quote}

These headers allow the quantifier order and exact coverage to be
checked without inferring them from theorem names. The recorded audit
contains no \nolinkurl{sorryAx} among the axioms of the examined
declarations. The supplementary report distinguishes that transitive
check from the textual search for \nolinkurl{native_decide} in project
sources and from a complete rebuild of dependencies.

\subsection{3. Dependencies, adapters and
attribution}\label{dependencies-adapters-and-attribution}

{\def\LTcaptype{none} % do not increment counter
\begin{longtable}[]{@{}
  >{\raggedright\arraybackslash}p{(\linewidth - 2\tabcolsep) * \real{0.5000}}
  >{\raggedright\arraybackslash}p{(\linewidth - 2\tabcolsep) * \real{0.5000}}@{}}
\toprule\noalign{}
\begin{minipage}[b]{\linewidth}\raggedright
Result in the text
\end{minipage} & \begin{minipage}[b]{\linewidth}\raggedright
Reference module
\end{minipage} \\
\midrule\noalign{}
\endhead
\bottomrule\noalign{}
\endlastfoot
\textbf{Inputs and construction} & \\
Attained rational optimum & \nolinkurl{PaperI.FiniteLPDuality},
\nolinkurl{CertifiedOptimumExistence} \\
Uniform mixed rounding, Theorem 3.1 & \nolinkurl{RC01Final},
\nolinkurl{MixedRoundingAdapter}; \nolinkurl{Nibble.*} infrastructure from
Paper III \\
Corollary 3.5, far regime for all graphs &
\nolinkurl{FarRegimeAllGraphs.farRegime_cliquePartition_allGraphs};
explicit version in \nolinkurl{E17.ExplicitFarAssembly} \\
Complete improved cleaning and explicit threshold &
\nolinkurl{E17.Cleanup}, \nolinkurl{E17.FarBudget},
\nolinkurl{E17.ExplicitFarRounding}, \nolinkurl{E17.ExplicitFarAssembly};
constants from \nolinkurl{E18} and \nolinkurl{E19} \\
Fractional nibble selection, Paper III & \nolinkurl{Nibble.FracNibbleLE},
\nolinkurl{Nibble.FracNibbleRepaired}; adapters
\nolinkurl{PaperIIISlackNibbleAdapter}, \nolinkurl{PaperIIINibbleAdapter} \\
Localization, root and critical window & \nolinkurl{NearH1Localization},
\nolinkurl{NearRootWindow}, \nolinkurl{RootRegularizationBridge} \\
Dichotomy and absorbable interface & \nolinkurl{HybridDichotomy},
\nolinkurl{SeparationAbsorptionRoute} \\
Accounting and near partition & \nolinkurl{RD09PhysicalLedger},
\nolinkurl{NearH1FinalAssembly}, \nolinkurl{NearH1LocalConstructor} \\
\textbf{Extremal conclusions} & \\
Sharp bound and eventual maximum & \nolinkurl{Erdos81Unconditional},
\nolinkurl{PaperTheorems} \\
Bounds for all orders & \nolinkurl{Erdos81AllOrders}; explicit bound in
\nolinkurl{ExplicitThreshold} \\
Fixed rooted defect and maximum, Theorem C &
\nolinkurl{DefectExplicitPublication}; explicit assembly
\nolinkurl{E34.L11Main}; lower witness \nolinkurl{DefectSharpPublication} \\
Necessity of order four, Proposition D.2 &
\nolinkurl{DefectSharpPublication.order_three_insufficient},
\nolinkurl{E11K10} \\
General budget and passage to \(\operatorname{cp}\) &
\nolinkurl{LossBudget}, \nolinkurl{SplitMixedGap} \\
Graph stability, partition stability and extremal classification &
\nolinkurl{IntegralStability}, \nolinkurl{RootPartitionStability},
\nolinkurl{ExtremalClassification} \\
Complete-split value and cores & \nolinkurl{SplitCompleteSharpValue},
\nolinkurl{SplitCompleteRigidity} \\
Partition defect and optimal coefficients &
\nolinkurl{SplitCompleteDefect}, \nolinkurl{LinearCoefficient},
\nolinkurl{SharpConstantOptimality} \\
Edge density in the critical branch & \nolinkurl{NearEdgeDensity},
\nolinkurl{NearCriticalDichotomy} \\
Equality of mixed optima, Appendix D.2 & \nolinkurl{SplitMixedGap},
\nolinkurl{SplitCompleteSharpValue} \\
\textbf{Complementary tools} & \\
Budgeted compatible absorption &
\nolinkurl{SpreadAbsorptionCompatibility} \\
Balanced reserve, Appendix B.2 & \nolinkurl{BalancedReserve},
\nolinkurl{Nibble.BeckFiala} from Paper III \\
Four-host completion, Appendix B.3 & \nolinkurl{FourHostClosure}; Vizing
and \nolinkurl{MultiHostTriangleLift} \\
Exact reserve and threshold sensitivity & \nolinkurl{ReserveIdentity},
\nolinkurl{NearThresholdSensitivity} \\
Reusable chordal structure & \nolinkurl{CliqueTree} \\
\end{longtable}
}

\textbf{Module map.} The prefixes \nolinkurl{PaperI}, \nolinkurl{Nibble} and
\nolinkurl{MixedRounding} indicate, respectively, the rational development
associated with Paper I, the nibble infrastructure of Paper III and the
neutral mixed-model library. The prefixes \nolinkurl{E17}, \nolinkurl{E18},
\nolinkurl{E19} and \nolinkurl{A4S1} identify the increment modules integrated
into the project; other unprefixed names belong to \nolinkurl{PaperIV};
\nolinkurl{MixedRoundingAdapter} adapts that namespace to the
\nolinkurl{MixedRounding} library. RC01 incorporates and adapts tools from
Paper III, but its two-quota mixed assembly is the one in Section 3 of
the present work. The table relates results to sources; the supplement
retains the exact dependency list.

Chordal structure and Dirac's theorem come from the Paper II
contribution, ported into \nolinkurl{ChordalStructure}. Szemerédi's
regularity lemma is reused from Mathlib; the construction and adapters
for clean profiles and mixed selection belong to Paper IV. This
provenance is distinct from a file's current location under
\nolinkurl{PaperIV}.

The configuration declares Mathlib as a third-party package and a
selection of Paper III \nolinkurl{Nibble} modules as a local dependency of
the series. \nolinkurl{ConeAudit} traverses the types and proof terms of
the selected declarations and checks that their transitive dependencies
contain no names under the root \nolinkurl{Erdos81}, used by {[}5{]}. Our
modules \nolinkurl{PaperIV.Erdos81Unconditional} and
\nolinkurl{PaperIV.Erdos81AllOrders} have namespace root \nolinkurl{PaperIV};
they are not the excluded root Erdos81 used by {[}5{]}. The absence of
these names is a reproducible audit result, not an inference from the
import list. Its scope is precise: it certifies that exclusion in the
examined snapshot; it would not by itself detect relocated and renamed
code or determine the historical provenance of each argument. Source
attribution is documented separately. The same namespace check cannot
yield a verdict on {[}15{]}. The extension uses the signed-star method
of {[}15, Lemma 3.1{]} and the removal method of {[}22{]}; proving their
implementations in this supplement does not make those methods new.

The public interface for Theorem C in this revision is
\nolinkurl{DefectExplicitPublication}: it uses
\nolinkurl{E34.theoremC_fully_explicit_final} and the unrestricted lower
witness already proved in \nolinkurl{DefectSharpPublication}. The older
existential assembly is not the proof selected for Theorem C. Its
localization uses induced removal of Alon--Shapira {[}23{]}; the
selected explicit assembly instead uses the polynomial chordal-removal
development of {[}22{]}. The additional \nolinkurl{FDCheck.ASCheck}
traverses proof terms and types to test exclusion of the named removal
results in these selected declarations. This is not exclusion of the
\nolinkurl{AlonShapira} namespace: shared definitions of chordality, edit
distance and induced copies remain in use.

There is a specific exception to this dependency comparison. The
unrestricted equality classification in Corollary 6.3, (6.18), is
exported through \nolinkurl{E32.cp_classification_of_theoremCPrimeC};
that declaration and \nolinkurl{E32.ref15Theorem11_of_theoremCPrimeC}
still use the historical assembly through \nolinkurl{E32.theoremC_at}.
Their proofs therefore include the induced-removal chain of {[}23{]}.
All those lemmas are proved in the frozen sources with standard axioms.
The exclusion check for the explicit Theorem C and the selected
stability declarations must not be extended to these classification
wrappers.

\subsubsection{3.1. Resources and the two gain
conventions}\label{resources-and-the-two-gain-conventions}

The representation refinement in the original proof map is verified in
\nolinkurl{MixedRoundingAdapter}. The adapter identifies a piece's edges
and the admissible-copy predicate, transports fractional packing
capacities, and proves equality of values via \nolinkurl{value_toFarFrac}.
For an integral packing, \nolinkurl{gain_toFarPacking} preserves gain.
Thus Theorem 3.1 can be used through this neutral presentation without
changing budget (1.3). This is not a new proof of rounding.

Two functions named \nolinkurl{gainOf} must be distinguished by their
domain of use. \nolinkurl{Model.gainOf} is two at cardinality three, five
at cardinality four, and zero at all other cardinalities.
\nolinkurl{FarRounding.gainOf} is \(\binom{|K|}{2}\mathbin{\dot-}1\), where
\(\dot-\) denotes truncated subtraction in the naturals. They agree at
cardinalities zero through four, but cease to agree from five onwards.
Mixed items have cardinality three or four, so the transport is
legitimate. For unrestricted partitions, such as those in Section 6.4,
one uses edge accounting or \nolinkurl{FarRounding.gainOf}; it is never
replaced by the bounded gain function from \nolinkurl{Model}. Since
partition pieces have at least two vertices, truncated subtraction there
agrees with ordinary subtraction.

\subsubsection{3.2. Construction and retained
accounts}\label{construction-and-retained-accounts}

The bounded-rank nibble is reused through
\nolinkurl{PaperIIISlackNibbleAdapter.boundedRankNibbleAt}. The two-quota
construction is in
\nolinkurl{MarkedQuotaSlackGate.slackMarkedQuotaNibbleAt_proved}; the
pairing of triangular supports and its physical realization are in
\nolinkurl{MarkedQuotaPairing} and \nolinkurl{JointTwoQuotaPhysical}. The
joint codegree scale corresponds to
\nolinkurl{RC01CleanedGate.cleanedPacking_joint_codegree_le_of_served_patterns}.
Profile mass is at most \(t^2\), not one: this is the physical
normalization of \nolinkurl{RC01PatternMassScale}.

For the mixed path, \nolinkurl{VertexCopyMonotone.two_mul_F4_le_add}
proves the two-direction inequality. \nolinkurl{VertexCopyGate} selects the
step and its finite potential;
\nolinkurl{GatedTerminalSplit.exists_split_terminal_symmetrizationPath}
supplies the complete sequence. Local accounting is constructed in
\nolinkurl{NearH1WindowAccounts} before being imported into
\nolinkurl{NearH1Localization}. This dependency direction matches the
logical order explained in Lemma 4.2.

Clique extraction and calibration are in \nolinkurl{ChordalCoreMissing} and
\nolinkurl{NearH1CalibratedRoot}. \nolinkurl{Regularization},
\nolinkurl{RegularizationBounds} and \nolinkurl{RegularizedRootGoal} prove the
root properties used by \nolinkurl{RootRegularizationBridge}. Finally,
\nolinkurl{NearH1PhaseI}, \nolinkurl{RD09FactorCandidateMoments} and
\nolinkurl{RD09FactorCandidateAverage} produce the two accounts;
\nolinkurl{NearH1FinalAssembly} combines them using the same first
assignment and its occupied links. The second phase is not applied to an
independent copy of the resources.

Theorem 5.0 has a strengthened form,
\nolinkurl{NearH1LocalConstructor.exists_near_partition_paid_by_root_sharp},
which adds \(|3|R|-n|\le n/50\) to the physical constructor; the
previous interface \nolinkurl{exists_near_partition_paid_by_root}
remains as a corollary. The strengthened reserve in (5.15a) is retained
for the same witness in
\nolinkurl{NearStructureWitness.accounts_count_le_improved};
\nolinkurl{IntegralStability.chordal_linear_stability_sixteen} and
\nolinkurl{RootPartitionStability.chordal_partition_stability} use it
without recombining roots or packings. The cyclic selection in phase I
is in \nolinkurl{RD09PaddedL1Mass}; the final evaluation of phase II, in
\nolinkurl{H1ImprovedConstants.l10_of_raw_rd09L2}. The comparator
regimes in (5.1a) are checked in
\nolinkurl{SplitComparatorResidual.residual_sq_le_of_near_split_universal}.
Reserve (D.8) is
\nolinkurl{ReserveIdentity.targetSize_eq_baseline_add_reserve}; the
improved contraction corresponds to
\nolinkurl{NearThresholdSensitivity.descent_threshold_maximal_barrier}
and \nolinkurl{near_threshold_binding}. The critical-geometry results are
in \nolinkurl{NearRootWindow}, \nolinkurl{NearEdgeDensity} and
\nolinkurl{NearCriticalDichotomy}.

The extension results have separate public entry points. Theorem C and
its maximum are
\nolinkurl{PaperIV.DefectExplicitPublication.rooted_defect_eventual} and
\nolinkurl{rooted_defect_maximum}; the historical
\nolinkurl{DefectSharpPublication} namespace retains the lower witness and
\nolinkurl{order_three_insufficient} for Proposition D.2.
\nolinkurl{E35.theoremC_tower} and its uniform variants give Corollary
6.6. The complete account (3.11) is
\nolinkurl{E17Bridge.improved_far_loss_explicit}. Corollary 3.5 with an
explicit threshold is
\nolinkurl{E17Bridge.ExplicitAssembly.farRegime_allGraphs_explicit}.
Finally, \nolinkurl{PaperIV.ExplicitThreshold.erdos81_sharp_explicit} and
\nolinkurl{erdos81_all_orders_bounded_additive} give the global result
with the tower of §6.3. These names identify declarations, not new
hypotheses.

\subsubsection{3.3. Scope relative to other formal
developments}\label{scope-relative-to-other-formal-developments}

The formal scope differs from that documented by {[}5{]} at the cited
commit. Its \nolinkurl{lean/FORMALIZATION_STATUS.md} file declares
verification conditional on \nolinkurl{ExternalInputs.Inputs}: Vizing,
Häggkvist--Janssen and a packing transfer that includes certificates of
rational optimality. These inputs are explicit parameters, not hidden
axioms. Here the inputs used by our chain, including Vizing, are
discharged; Häggkvist--Janssen is not used. The comparison concerns what
each Lean development checks: it does not assert that the mathematical
proof of {[}5{]} is conditional on conjectures.

Reusable contributions have their own references. Sum-zero triangle
packing corresponds to Lean Pool PR \#348 {[}12{]},
\url{https://github.com/Vilin97/lean-pool/pull/348}, merged at commit
\nolinkurl{540d8e3}. Minimum-degree matchings and their weighted selection
correspond to PR \#420 {[}9{]},
\url{https://github.com/Vilin97/lean-pool/pull/420}, merged at
\nolinkurl{d1de6d2}. The series sources retain the antecedents and
adaptations of Vizing and Beck--Fiala. Attribution of formalizations is
distinguished from authorship of the classical theorems.

Lean {[}10{]} and Mathlib {[}11{]} provide the verification environment.
Code auditing does not replace human review of definitions or of the
correspondence between a formal statement and its expression in the
manuscript. Experiments guided the search and detected false assertions;
they are not premises of universality.

The formal constructor used here invokes neither the results of {[}15{]}
nor the historical modules of the replaced construction. The list bound
and first-phase capacity window have the same form as conditions in
{[}15, Lemma 4.1 and (4.4){]}. Here they are justified by dyadic
hosting, averaging over cyclic shifts, and balanced absorption. This
comparison identifies the mechanisms and their shared conditions; it
claims neither priority for those conditions nor historical independence
on the basis of the dependency check in §3 of this companion.

\subsection{4. Precision notes for the
manuscript}\label{precision-notes-for-the-manuscript}

The following distinctions survive the removal of implementation names
from the exposition. They prevent a mathematical statement from being
identified with a stronger or differently quantified Lean type.

\begin{itemize}
\tightlist
\item
  The construction in §6.1 uses triangles and edges; the audited public
  interface supplies a partition of order at most four and an
  unrestricted lower bound.
\item
  The quadratic and linear optimality declarations use rational
  coefficients. The manuscript explains the passage to real coefficients
  by rational approximation.
\item
  \nolinkurl{recover_clique_from_edit} takes rational \(u\ge0\); the
  written proof also applies to real \(u\ge0\).
\item
  Proposition G.1 uses a rational precision parameter in the Lean
  declaration and real exponential comparisons.
\item
  The width formulation in Appendix G.1 uses clique trees with bags of
  bounded cardinality; it does not rely on a general Mathlib treewidth
  definition.
\item
  No chordal restriction of the cleanup-to-rounding implication is
  formalized by Appendix G.2. The universal obstruction is not a lower
  bound for every chordal rounding method.
\end{itemize}

\subsubsection{4.1. An existing numerical premise made
explicit}\label{an-existing-numerical-premise-made-explicit}

The old prose in F.5 referred to \nolinkurl{thm2_ineq} by name. The new
prose prints its existing conclusion, with \(m=m_{11}(\varepsilon,K)\),
\(B=B_{11}(\varepsilon,K)\) and \(\delta_{11}\) as in (F.11). The source
is \nolinkurl{E34/L11Params.lean}, lines 289--292. Its assumptions are
\(0<\varepsilon<1\) and \(1\le K\); the preceding paragraph of F.5
supplies the stronger \(\varepsilon<1/2\) and \(K\ge1\). This is an
expansion of a cited premise, not a new estimate or proof.

\subsubsection{4.2. Source map for relocated inline
references}\label{source-map-for-relocated-inline-references}

{\def\LTcaptype{none} % do not increment counter
\begin{longtable}[]{@{}
  >{\raggedright\arraybackslash}p{(\linewidth - 2\tabcolsep) * \real{0.5000}}
  >{\raggedright\arraybackslash}p{(\linewidth - 2\tabcolsep) * \real{0.5000}}@{}}
\toprule\noalign{}
\begin{minipage}[b]{\linewidth}\raggedright
Manuscript location
\end{minipage} & \begin{minipage}[b]{\linewidth}\raggedright
Formal declarations or modules
\end{minipage} \\
\midrule\noalign{}
\endhead
\bottomrule\noalign{}
\endlastfoot
1.1. The common condition: a loss budget &
\nolinkurl{LossBudget.budget_iff} \\
1.2. Structural dichotomy and resolution mechanisms &
\nolinkurl{NearStructureWitness} \\
5.3. The three accounts that pay for the construction &
\nolinkurl{isSpokeCompatible_surviving_canonicalBad},
\nolinkurl{RD09SpokeCompatibility.card_union_phases} \\
6.2. The eventual maximum &
\nolinkurl{SplitCompleteRigidity.optimal_cores} \\
6.2. The eventual maximum &
\nolinkurl{PaperTheorems.splitGraph_isChordal} \\
6.3. All orders and optimality of the linear term &
\nolinkurl{SharpConstantOptimality.erdos81_quadratic_constant_optimal},
\nolinkurl{erdos81_quadratic_constant_isLeast} \\
6.3. All orders and optimality of the linear term &
\nolinkurl{LinearCoefficient.linear_coefficient_optimal} \\
6.3. All orders and optimality of the linear term &
\nolinkurl{LossBudget.erdos81_cp_form_all_orders},
\nolinkurl{erdos81_cp_form} \\
6.6. Fixed-defect stability on one root & \nolinkurl{NormalizedDeficit},
\nolinkurl{ResidualExcess}, \nolinkurl{NormalizedCliqueCore},
\nolinkurl{RootPartition}, \nolinkurl{NormalizedStability} \\
6.7. Uniform bounds and sublinear defect &
\nolinkurl{E35.Fexp_le_tower} \\
B.1. Budgeted compatible absorption &
\nolinkurl{SpreadAbsorptionCompatibility.exists_budgeted_compatible_spread_absorber} \\
B.2. Balanced reserves &
\nolinkurl{BalancedReserve.exists_balanced_reserve} \\
B.3. Completion with four hosts &
\nolinkurl{FourHostClosure.exists_fourHost_packing} \\
C.2. What is discarded and why the threshold is uniform &
\nolinkurl{RC01CleanFiber}, \nolinkurl{RC01DeviationCleanup} \\
C.3. An explicit schedule and its propagation &
\nolinkurl{SzemerediRegularity.initialBound} \\
C.3. An explicit schedule and its propagation &
\nolinkurl{E18.NibbleSchedule}, \nolinkurl{NibbleOracle},
\nolinkurl{NibbleChain} \\
C.3. An explicit schedule and its propagation &
\nolinkurl{E19.ScheduleBounds}, \nolinkurl{E19.ChainBounds},
\nolinkurl{E19.GateBounds} \\
C.3. An explicit schedule and its propagation & \nolinkurl{E18.Numeric},
\nolinkurl{E19.Main} \\
C.3. An explicit schedule and its propagation &
\nolinkurl{E35.NfarE_le_tower_poly} \\
D.1. The complete graph at every order &
\nolinkurl{ThreeRegime.CompleteStateAllOrders} \\
D.1. The complete graph at every order & \nolinkurl{HalvingFrames} \\
D.1. The complete graph at every order & \nolinkurl{exists_packing},
\nolinkurl{complete_state_closes} \\
D.2. Equality of mixed optima in complete-split graphs &
\nolinkurl{SplitMixedGap.mixed_gap_zero},
\nolinkurl{SplitCompleteSharpValue} \\
E.1. Constructor and accounting & \nolinkurl{IndepAllBounds},
\nolinkurl{IndepAllParams} \\
E.1. Constructor and accounting & \nolinkurl{IndepAllPeel},
\nolinkurl{IndepAllObstr} \\
E.1. Constructor and accounting & \nolinkurl{IndepAllBudget} \\
E.1. Constructor and accounting & \nolinkurl{IndepAllMain},
\nolinkurl{ConstructorBudget} \\
E.2. From the terminal to the eventual theorem &
\nolinkurl{E34.theoremC_fully_explicit_final} \\
E.4. Retained accounts for same-root stability &
\nolinkurl{StrictParams} \\
E.4. Retained accounts for same-root stability &
\nolinkurl{StrictBudget} \\
E.4. Retained accounts for same-root stability &
\nolinkurl{ResidualExcess} \\
E.4. Retained accounts for same-root stability &
\nolinkurl{AllInput.no_core_pairs} \\
E.4. Retained accounts for same-root stability &
\nolinkurl{CoreCliqueAlternative} \\
E.4. Retained accounts for same-root stability &
\nolinkurl{TemplateDistance}, \nolinkurl{RootPartition},
\nolinkurl{CliqueShiftArithmetic} \\
F.1. Explicit fixed-defect parameters & \nolinkurl{ExplicitFixedStability},
\nolinkurl{E33.Assembly}, \nolinkurl{E33.LocExplicit}, \nolinkurl{E34.L11Main} \\
F.1. Explicit fixed-defect parameters &
\nolinkurl{fixed_defect_stability_explicit} \\
F.2. Fractional star accounting and small orders &
\nolinkurl{RefinedFractionalBound}, \nolinkurl{UniformFractionalBound},
\nolinkurl{SmallOrderFractional} \\
F.3. Uniform localization through a chordal edit &
\nolinkurl{SublinearEdit} \\
F.3. Uniform localization through a chordal edit &
\nolinkurl{exists_clique_additive_defect},
\nolinkurl{recover_clique_from_edit} \\
F.3. Uniform localization through a chordal edit &
\nolinkurl{SublinearStabilityTransfer} \\
F.5. Auxiliary lemmas and adaptations of chordal removal &
\nolinkurl{E34.indCopies_cycle_le} \\
F.5. Auxiliary lemmas and adaptations of chordal removal &
\nolinkurl{E34.coBipartite_edit_le} \\
F.5. Auxiliary lemmas and adaptations of chordal removal &
\nolinkurl{E34.exists_large_clique} \\
F.5. Auxiliary lemmas and adaptations of chordal removal &
\nolinkurl{E34.lemma3} \\
F.5. Auxiliary lemmas and adaptations of chordal removal &
\nolinkurl{E34.transfer} \\
F.5. Auxiliary lemmas and adaptations of chordal removal &
\nolinkurl{E34.lemma11_V} \\
F.5. Auxiliary lemmas and adaptations of chordal removal &
\nolinkurl{properOn_of_pinned} \\
F.5. Auxiliary lemmas and adaptations of chordal removal &
\nolinkurl{thm2_ineq} \\
F.5. Auxiliary lemmas and adaptations of chordal removal &
\nolinkurl{claim5_count}, \nolinkurl{block_ineq} \\
F.5. Auxiliary lemmas and adaptations of chordal removal &
\nolinkurl{l11_caseA}, \nolinkurl{l11_caseB1}, \nolinkurl{l11_caseB2},
\nolinkurl{E34.L11Main} \\
F.5. Auxiliary lemmas and adaptations of chordal removal &
\nolinkurl{E35.theoremC_tower}, \nolinkurl{theoremC_tower_uniform},
\nolinkurl{theoremC_tower_uniform_sMax},
\nolinkurl{NfarE_le_tower_poly},
\nolinkurl{PaperIV.DefectExplicitPublication}, \nolinkurl{FDCheck.ASCheck},
\nolinkurl{PaperIV.OptimalTemplateObstruction.shifted_template_witness},
\nolinkurl{shifted_template_sqrt_witness},
\nolinkurl{no_linear_optimal_template_bound},
\nolinkurl{optimal_family_nonempty} \\
G.1. The all-order bound and packing gaps & \nolinkurl{BoundedCliqueGap} \\
G.1. The all-order bound and packing gaps &
\nolinkurl{BoundedCliqueGap.chordal_gap_linear_cliqueFree} \\
G.1. The all-order bound and packing gaps &
\nolinkurl{exists_cliqueTree_width_le_iff_cliqueFree} \\
G.2. A quantitative obstruction to codegree cleanup &
\nolinkurl{CleanupAtWith} \\
G.2. A quantitative obstruction to codegree cleanup &
\nolinkurl{threshold_gt_exp}, \nolinkurl{threshold_gt_exp_seventy_two},
\nolinkurl{FarExploration.CleanupRigidVerdict} \\
G.2. A quantitative obstruction to codegree cleanup &
\nolinkurl{threshold_superpolynomial} \\
G.2. A quantitative obstruction to codegree cleanup &
\nolinkurl{CodegreeCleanupAt}, \nolinkurl{UniformRoundingTarget} \\
\end{longtable}
}

\subsection{5. Reproduction and separate
annex}\label{reproduction-and-separate-annex}

Run the following from the extracted, unchanged source archive with the
toolchain pinned by \nolinkurl{lean-toolchain} and dependencies pinned by
\nolinkurl{lake-manifest.json}. The command is documented for reproduction;
it was not executed for this editorial revision.

\begin{verbatim}
python tools/audit_publication.py
\end{verbatim}

It reads the 19 targets in \nolinkurl{FREEZE_SCOPE.json}, runs them
serially and keeps individual logs. Reimporting a compiled audit module
is not equivalent to rerunning its commands. The recorded run is under
\nolinkurl{03_reproducibility/full_rebuild_v12_20260929_resume/}, with
provenance links to the first segment. \nolinkurl{ReleaseExportCheck} uses
only \texttt{import\ PaperIV}; the named audits inspect their own lists
of types and proof terms.

{\def\LTcaptype{none} % do not increment counter
\begin{longtable}[]{@{}l@{}}
\toprule\noalign{}
Target \\
\midrule\noalign{}
\endhead
\bottomrule\noalign{}
\endlastfoot
\nolinkurl{PaperIV} \\
\nolinkurl{PaperIV.Audit} \\
\nolinkurl{PaperIV.ConeAudit} \\
\nolinkurl{PaperIV.PublicationIncrementAudit} \\
\nolinkurl{A4S1.IndepAllAudit} \\
\nolinkurl{ThreeRegime.AxiomCheck} \\
\nolinkurl{FarExploration.AxiomCheck} \\
\nolinkurl{PaperIV.CliqueTreeAudit} \\
\nolinkurl{E17.ImprovedFarRoundingAudit} \\
\nolinkurl{E17.ExplicitFarAudit} \\
\nolinkurl{PaperIV.ExplicitThresholdAudit} \\
\nolinkurl{ReleaseExportCheck} \\
\nolinkurl{PaperIVEditorial.ConstructorBudgetAudit} \\
\nolinkurl{ResearchAudit} \\
\nolinkurl{FDCheck.FinalAudit} \\
\nolinkurl{E35.Audit} \\
\nolinkurl{FDCheck.ASCheck} \\
\nolinkurl{E34.Audit} \\
\nolinkurl{PaperIV.OptimalTemplateObstructionAudit} \\
\end{longtable}
}

The literal source manifest is
\nolinkurl{03_reproducibility/build_piv-v12-fb459343d234/SOURCE_MANIFEST.json}.
The export log records 224 literal \nolinkurl{#check} commands. The
constructor audit covers \nolinkurl{PaperIV.Editorial.ConstructorBudget},
which keeps both conclusions for one partition. The old extraction
workspaces and superseded source cuts remain historical records in the
series repository; they are not extra modules of the freeze described
here.

\nolinkurl{BoundedCliqueGap} is a separate source annex supporting Appendix
G.1; it is not an input to the main extremal proof. Its source archive
SHA-256 is
\nolinkurl{2847a422865e06880d457ea806aae5b9f749d22359eff325349c09c01cb4837d}.
Its 38 source modules and 50-module build closure belong to its distinct
recorded build, not to the count of 607 above.

\subsection{6. Written-proof review and
limits}\label{written-proof-review-and-limits}

\textbf{Scope of the written proof.} The table below records seven
derivations examined in the consolidated external review. Its
assessment, ACCEPTABLE\_SUMMARY for each row, retains the stated
expository limit concerning the cited second-moment estimates in C.2. A
successful build checks formal derivations, not completeness of
exposition; the written-proof assessment remains distinct from
compilation and from human peer review. The audit history and detailed
assessments are retained with the supplementary evidence.

{\def\LTcaptype{none} % do not increment counter
\begin{longtable}[]{@{}
  >{\raggedright\arraybackslash}p{(\linewidth - 2\tabcolsep) * \real{0.5000}}
  >{\raggedright\arraybackslash}p{(\linewidth - 2\tabcolsep) * \real{0.5000}}@{}}
\toprule\noalign{}
\begin{minipage}[b]{\linewidth}\raggedright
Location
\end{minipage} & \begin{minipage}[b]{\linewidth}\raggedright
Derivation to review, including the expanded accounts
\end{minipage} \\
\midrule\noalign{}
\endhead
\bottomrule\noalign{}
\endlastfoot
§5.1, (5.4) and (5.7) & Regularization mass, degree, palette and width
estimates \\
C.2, (3.8) and (3.10) & Discard counts and retention under the partition
conventions \\
Lemma 3.2 and C.1--C.2 & Fibre lower bounds from regular pairs and their
normalization \\
C.3, Table C.1 & Numerical domination of the nibble constants and tower
iterations \\
E.1 & Normalization bounds preceding the terminal construction \\
E.1, (E.2b) and (E.2d) & Terminal estimates and the three-case budget \\
D.1 & Bose/Skolem frame constructions underlying the residue table \\
\end{longtable}
}

Experimental checks of small orders are not used as premises. Their
documentation, including completeness and certification requirements,
belongs to the separate research record.

The clique-tree library retains counterexamples to formulations missing
the appropriate hypotheses: a leaf alone is not equivalent to a unique
simplicial vertex, and separator statements must account for repeated
bags and degenerate cases. These results delimit the library's use; they
add no assumptions to the main theorem.

The editorial ledger retains the exact source paragraphs and every
replacement. Automated checks compare displayed equations, labels and
result headings; hypotheses and proof dependencies are also reviewed in
the source-to-manuscript comparison. These checks do not automatically
certify semantic equivalence of all prose. The AI/tools disclosure
retains its attribution and responsibility statements; the superseded
draft-status sentence has been removed.

\subsection{Use of artificial intelligence and computational
tools}\label{use-of-artificial-intelligence-and-computational-tools}

Claude, by Anthropic, and ChatGPT/Codex, by OpenAI, were used to explore
and check arguments and to prepare the manuscript. Aristotle, by
Harmonic, contributed to exploring approaches, searching for
counterexamples, developing proofs, and formalizing and reviewing them
in Lean; its role was not limited to translating finished proofs.

Certo {[}13{]}, \url{https://github.com/jtraverso/certo-math}, and
Jacobian {[}14{]}, \url{https://github.com/morluto/jacobian}, were the
computational tools used to explore instances, check identities, and
assess candidates. Jacobian is also credited by the project behind
{[}5{]}; its use here is acknowledged separately from the logical
dependencies of our formal proof. Computations were reviewed within
their stated scope; a finite search does not replace a universal proof.

The author remains responsible for the arguments, citations, code, and
presentation. No AI system is listed as an author.

\subsection{References}\label{references}

{[}1{]} P. Erdős, E. T. Ordman, and Y. Zalcstein, ``Clique partitions of
chordal graphs,'' \emph{Combinatorics, Probability and Computing}
\textbf{2} (1993), 409--415.

{[}2{]} J. P. Traverso Gianini, \emph{Affine Profile Reduction for
Fractional Triangle Packings in Split Graphs}, Paper I, preprint v1.3,
22 August 2026. Spanish and English editions in the joint v3 deposit, 23
August 2026. Version DOI consulted:
\url{https://doi.org/10.5281/zenodo.22064657}; concept DOI:
\url{https://doi.org/10.5281/zenodo.21273143}. Supplementary material:
\url{https://github.com/jtraverso/erdos-81-chordal-clique-partitions/tree/main/preprints/PAPER_I}.

{[}3{]} J. P. Traverso Gianini, \emph{Complete-Split Extremizers for a
Fractional Triangle-Cover Functional on Chordal Graphs}, Paper II,
preprint v1.2, 22 August 2026. Spanish and English editions in the joint
v3 deposit, 23 August 2026. Version DOI consulted:
\url{https://doi.org/10.5281/zenodo.22064657}; concept DOI:
\url{https://doi.org/10.5281/zenodo.21273143}. Supplementary material:
\url{https://github.com/jtraverso/erdos-81-chordal-clique-partitions/tree/main/preprints/PAPER_II}.

{[}4{]} J. P. Traverso Gianini, \emph{Linear-Error Clique Partitions of
Split Graphs via Structured Triangle Packing}, Paper III, preprint v1.5,
23 August 2026. Spanish and English editions in the joint v3 deposit of
that date. Version DOI consulted:
\url{https://doi.org/10.5281/zenodo.22064657}; concept DOI:
\url{https://doi.org/10.5281/zenodo.21273143}. Supplementary material:
\url{https://github.com/jtraverso/erdos-81-chordal-clique-partitions/tree/main/preprints/PAPER_III}.

{[}5{]} Anonymous, \emph{Clique partitions of chordal graphs with linear
error}, preprint, 8 September 2026. Version consulted on 20 September
2026, commit \nolinkurl{cbde8a0a0563372b23b1b39a44180d2c0fb02f44}.
Manuscript:
\url{https://github.com/N0zoM1z0/erdos-81/blob/cbde8a0a0563372b23b1b39a44180d2c0fb02f44/manuscript/main.pdf}.
The attribution printed in the document is retained. The README at that
same commit credits Morluto for the mathematical proof, N0zoM1z0 for the
manuscript and Lean implementation, and Jacobian at Preference Labs as
part of the project:
\url{https://github.com/N0zoM1z0/erdos-81/blob/cbde8a0a0563372b23b1b39a44180d2c0fb02f44/README.md}.
These repository credits are distinguished from the signature printed on
the manuscript.

{[}6{]} P. E. Haxell and V. Rödl, ``Integer and fractional packings in
dense graphs,'' \emph{Combinatorica} \textbf{21} (2001), 13--38.

{[}7{]} R. Yuster, ``Integer and fractional packing of families of
graphs,'' \emph{Random Structures \& Algorithms} \textbf{26} (2005),
110--118.

{[}8{]} T. F. Bloom, ``Erdős Problem \#81,'' \emph{Erdős Problems}.
Available at \url{https://www.erdosproblems.com/81}. Cited to identify
the problem; this article does not certify the current review status of
published proposals.

{[}9{]} J. P. Traverso Gianini and contributors to the assisted
formalization, ``Minimum-degree and spread matching theorems,''
contribution to \emph{lean-pool}, PR \#420, 2026. Available at
\url{https://github.com/Vilin97/lean-pool/pull/420}; merge commit
\nolinkurl{d1de6d2}, 13 September 2026. The history and file headers retain
the attribution of the Aristotle-assisted contribution.

{[}10{]} L. de Moura and S. Ullrich, ``The Lean 4 Theorem Prover and
Programming Language,'' in \emph{Automated Deduction -- CADE 28}, LNCS
\textbf{12699}, Springer, 2021, 625--635.

{[}11{]} The mathlib Community, ``The Lean Mathematical Library,'' in
\emph{CPP 2020}, ACM, 2020, 367--381.

{[}12{]} J. P. Traverso Gianini, ``Sum-zero triangle packing
formalization,'' Paper III contribution to \emph{lean-pool}, PR \#348,
2026. Available at \url{https://github.com/Vilin97/lean-pool/pull/348};
merge commit \nolinkurl{540d8e3}, 25 August 2026. See the formalization and
assistance credits in the sources.

{[}13{]} J. P. Traverso Gianini, \emph{Certo: herramientas
computacionales con certificados verificables}. Available at
\url{https://github.com/jtraverso/certo-math}. The version used in the
experiments must be fixed together with its certificates in the
supplement.

{[}14{]} Morluto and contributors, \emph{Jacobian: herramientas
matemáticas componibles}. Available at
\url{https://github.com/morluto/jacobian}. Component attribution and
versions are retained in the repository.

{[}15{]} O. Okechukwu, \emph{Clique partitions and bounded simplicial
defect}, arXiv:2609.20871v1, 15 September 2026. Available at
\url{https://arxiv.org/abs/2609.20871v1}. Accessed 27 September 2026.

{[}16{]} C. Henderson, H. Koerts, E. Roberge, S. Spirkl, and R. Whitman,
\emph{Clique Partitions of Split Graphs}, project announced as in
preparation on R. Whitman's research page:
\url{https://sites.google.com/view/rebeccawhitman/research}. Accessed 27
September 2026; the page consulted links no manuscript against which to
compare its results.

{[}17{]} F. A. Behrend, ``On Sets of Integers Which Contain No Three
Terms in Arithmetical Progression,'' \emph{Proceedings of the National
Academy of Sciences} \textbf{32} (1946), 331--332.
\url{https://doi.org/10.1073/pnas.32.12.331}. The constant used in
Appendix G.2 corresponds to \nolinkurl{Behrend.roth_lower_bound} in
Mathlib v4.28.0.

{[}18{]} I. Z. Ruzsa and E. Szemerédi, ``Triple systems with no six
points carrying three triangles,'' \emph{Combinatorics}, vol.~II,
Colloquia Mathematica Societatis János Bolyai \textbf{18},
North-Holland, 1978, 939--945. Author's bibliographic record:
\url{https://www.renyi.hu/~szemered/pub.html}.

{[}19{]} R. C. Bose, ``On the construction of balanced incomplete block
designs,'' \emph{Annals of Eugenics} \textbf{9} (1939), 353--399.
\url{https://doi.org/10.1111/j.1469-1809.1939.tb02219.x}.

{[}20{]} T. Skolem, ``Some Remarks on the Triple Systems of Steiner,''
\emph{Mathematica Scandinavica} \textbf{6} (1958), 273--280.
\url{https://tidsskrift.dk/math/article/view/10551}.

{[}21{]} R. Cipollini, partial proof announcement for Erdős Problem
\#81, 10 August 2026. Author's summary:
\url{https://www.erdosproblems.com/forum/thread/81/proof-claims\#proof-claim-201};
linked manuscript: \url{https://www.overleaf.com/read/thjptfhgnmxc}.
Cited for the announced subquadratic-error bound, not as an
independently audited input.

{[}22{]} R. de Joannis de Verclos, \emph{Chordal graphs are easily
testable}, arXiv:1902.06135v1, 16 February 2019. Available at
\url{https://arxiv.org/abs/1902.06135v1}. The removal argument
underlying Appendix F.3 is attributed to this work; its implementation
and adaptations are documented in Appendix F.5.

{[}23{]} N. Alon and A. Shapira, \emph{A characterization of the
(natural) graph properties testable with one-sided error}, SIAM Journal
on Computing \textbf{37} (2008), 1703--1727. DOI:
\url{https://doi.org/10.1137/06064888X}. Used by the historical
existential assembly and by the unrestricted-classification wrappers
identified in §3 of this companion, not as the removal theorem selected
for the explicit public Theorem C.

\end{document}
